\question
В древнем королевстве Алгоритмии был создан магический артефакт, охраняющий вход в Сокровищницу Знаний. Артефакт активируется и открывает проход только при одновременном нажатии нескольких магических рун, между которыми существует сложная логическая связь.

Юный маг-стажер Элгор пытается записать логику работы артефакта в виде булевой функции, чтобы создать для него ключ. Он выделил четыре критически важные руны и обозначил их переменными:

\begin{itemize}
\item W – \textit{Руна Воды} активирована (W = 1)
\item F – \textit{Руна Огня} активирована (F = 1)
\item E – \textit{Руна Земли} активирована (E = 1)
\item A – \textit{Руна Воздуха} активирована (A = 1)
\end{itemize}

Элгор расшифровал древние скрижали и выяснил, что артефакт срабатывает по следующему правилу:

«Артефакт откроется, если Руна Воздуха активна, и при этом либо Руна Воды, либо Руна Огня неактивна. Однако если активны и Руна Воды, и Руна Огня, то артефакт откроется, только если неактивна Руна Земли».

Необходимо:
\begin{enumerate}
    \item Записать логическую функцию f(W, F, E, A);
    \item Составить СДНФ;
    \item Составить полином Жегалкина двумя способами;
    \item Сделать вывод о совпадении/несовпадении полиномов.
\end{enumerate}

  ---------------

Автор -- Полина Корень